% Part: model-theory
% Chapter: lindstrom
% Section: ls-property

\documentclass[../../../include/open-logic-section]{subfiles}

\begin{document}

\olfileid{mod}{lin}{lsp}

\olsection{خاصيتا التراص ولوفنهايم--سكولم}

وللتراص ولوفنهايم--سكولم نظيران في المنطقات المجردة، يحصلان بتعميم
التعريفين المعروفين على الوجه الآتي.

\begin{defn}
نثبت للمنطق المجرد $\tuple{\OLId{latin-looped}{L},\models_\OLId{latin-looped}{L}}$ \emph{خاصية التراص} إذا قبلت كل
مجموعة~$\OLId{greek-variable}{Gamma}$ من !!{sentence}s~$\OLId{latin-looped}{L}(\Lang{L})$ الإرضاء، متى قبلته
كل مجموعة $\OLId{greek-variable}{Gamma}_\OLMathNumeral{0} \subseteq \OLId{greek-variable}{Gamma}$ منتهية.
\end{defn}

\begin{defn}
للمنطق $\tuple{\OLId{latin-looped}{L},\models_\OLId{latin-looped}{L}}$ \emph{خاصية لوفنهايم--سكولم النزولية} إذا
كان لكل $\OLId{greek-variable}{Gamma}$ قابلة للإرضاء نموذج !!a{enumerable}.
\end{defn}


أما التشاكل الجزئي المعرّف في \olref[bas][pis]{defn:partialisom} فأمر «جبري» محض:
لا يرجع تعريفه إلى !!{sentence}s اللغة، وإنما يرجع إلى ما تعطيه
الـ!!{language}~$\Lang{L}$ من ثوابت للـ!!{structure}s. فلا يمتنع استعماله
في المنطقات المجردة. وقد علمنا، في منطق الرتبة الأولى، من
\olref[bas][pis]{thm:p-isom2} أن !!{structure}s اثنتين إذا تشاكلتا جزئيًا
تكافأتا أوليًا. لكن ذلك البرهان بعينه لا يجري في المنطقات المجردة؛ إذ قد لا يتأتى
الاستقراء على الـ!!{formula}s حين نأخذ $!E \in \OLId{latin-looped}{L}(\Lang{L})$
كيفما اتفق. ومع هذا تبقى النتيجة صحيحة متى توافرت خاصية لوفنهايم--سكولم.

\begin{thm}
\ollabel{thm:abstract-p-isom}
إذا كان $\tuple{\OLId{latin-looped}{L},\models_\OLId{latin-looped}{L}}$ منطقًا سويًا تتحقق فيه خاصية لوفنهايم--سكولم،
فإن كل !!{structure}s اثنتين متشاكلتين جزئيًا تكونان متكافئتين أوليًا في
$\tuple{\OLId{latin-looped}{L},\models_\OLId{latin-looped}{L}}$.
\end{thm}

\begin{proof}
لنفرض، طلبًا للخلف، أن $\Struct{M} \simeq_\OLId{latin-ordinary}{p} \Struct{N}$، وأن ثمة $!E$ تحقق
$\Struct{M} \models_\OLId{latin-looped}{L} !E$ و$\Struct{N} \not\models_\OLId{latin-looped}{L} !E$. وتجيز لنا خاصية
التشاكل افتراض أن $\Domain{\OLId{latin-looped}{M}}$ و$\Domain{\OLId{latin-looped}{N}}$ متباينان، كما تجيز خاصية
التوسع افتراض أن $!E \in \OLId{latin-looped}{L}(\Lang{L})$ للـ!!{language} المنتهية
$\Lang{L}$. ولتكن $\mathcal{I}$ مجموعة من التشاكلات الجزئية بين
$\Struct{M}$ و$\Struct{N}$، ولا يضر بعموم الحجة أن نشترط أيضًا أنه إذا
كان $\OLId{latin-ordinary}{p} \in \mathcal{I}$ و$\OLId{latin-ordinary}{q} \subseteq \OLId{latin-ordinary}{p}$، فإن $\OLId{latin-ordinary}{q} \in \mathcal{I}$.

$\Domain{\OLId{latin-looped}{M}}^{<\omega}$ مجموعة المتتاليات المنتهية من !!{element}s
$\Domain{\OLId{latin-looped}{M}}$. ولتكن $\OLId{latin-looped}{S}$ العلاقة الثلاثية على $\Domain{\OLId{latin-looped}{M}}^{<\omega}$
الدالة على وصل المتتاليات؛ ومعنى ذلك أنه إذا كانت
$\mathbf{a},\mathbf{b},\mathbf{c} \in \Domain{\OLId{latin-looped}{M}}^{<\omega}$، فإن
$\OLId{latin-looped}{S}(\mathbf{a},\mathbf{b},\mathbf{c})$ إذا وفقط إذا كانت $\mathbf{c}$
وصل $\mathbf{a}$ و$\mathbf{b}$. ولتكن $\OLId{latin-looped}{T}$ العلاقة الثلاثية التي يكون
فيها $\OLId{latin-looped}{T}(\mathbf{a},\OLId{latin-ordinary}{b},\mathbf{c})$، لـ$\OLId{latin-ordinary}{b} \in \Domain{\OLId{latin-looped}{M}}$ و
$\mathbf{a},\mathbf{c} \in \Domain{\OLId{latin-looped}{M}}^{<\omega}$، إذا وفقط إذا كان
$\mathbf{a}=\OLId{latin-ordinary}{a}_\OLMathNumeral{1},\dots,\OLId{latin-ordinary}{a}_\OLId{latin-ordinary}{n}$ و$\mathbf{c}=\OLId{latin-ordinary}{a}_\OLMathNumeral{1},\dots,\OLId{latin-ordinary}{a}_\OLId{latin-ordinary}{n},\OLId{latin-ordinary}{b}$. اختر
!!{predicate}s جديدين $\OLId{latin-looped}{P}$ و$\OLId{latin-looped}{Q}$ ذوي ثلاثة مواضع، وكوّن
الـ!!{structure}~$\Struct{M}^*$ التي كونها
$\Domain{\OLId{latin-looped}{M}} \cup \Domain{\OLId{latin-looped}{M}}^{<\omega}$، والتي تحتوي $\Struct{M}$
بنية جزئية وتفسر $\OLId{latin-looped}{P}$ و$\OLId{latin-looped}{Q}$ بعلاقتي الوصل $\OLId{latin-looped}{S}$ و$\OLId{latin-looped}{T}$ (فتكون
$\Struct{M}^*$ في الـ!!{language}~$\Lang{L} \cup \{\OLId{latin-looped}{P},\OLId{latin-looped}{Q}\}$).

عرّف $\Domain{\OLId{latin-looped}{N}}^{<\omega}$ و$\OLId{latin-looped}{S}'$ و$\OLId{latin-looped}{T}'$ و$\OLId{latin-looped}{P}'$ و$\OLId{latin-looped}{Q}'$ و$\Struct{N}^*$
على نحو مماثل. وحيث فرضنا أن $\Struct{M} \simeq_\OLId{latin-ordinary}{p} \Struct{N}$، وجدت
علاقة~$\OLId{latin-looped}{I}$ بين $\Domain{\OLId{latin-looped}{M}}^{<\omega}$ و$\Domain{\OLId{latin-looped}{N}}^{<\omega}$ بحيث
$\OLId{latin-looped}{I}(\mathbf{a},\mathbf{b})$ إذا وفقط إذا كانت $\mathbf{a}$ و
$\mathbf{b}$ متشاكلتين وتحققان شرط الذهاب والإياب من
\olref[bas][pis]{defn:partialisom}. ثم نبني $\Struct{U}$، وهي
الـ!!{structure} التي !!{domain}ها اتحاد !!{domain}s~$\Struct{M}^*$
و$\Struct{N}^*$، وتحتوي $\Struct{M}^*$ و$\Struct{N}^*$ بوصفهما
!!{structure}s جزئية، في الـ!!{language} ذات !!{predicate} ثنائي
إضافي~$\OLId{latin-looped}{R}$ يفسر بالعلاقة~$\OLId{latin-looped}{I}$ و!!{predicate}s تدل على الـ!!{domain}s
$\Domain{\OLId{latin-looped}{M}}^*$ و$\Domain{\OLId{latin-looped}{N}}^*$.

\begin{figure}[h]
  \centering
  \begin{tikzpicture}[node distance=2cm, auto, thick, >=stealth']
    \draw [rounded corners] (0,0) -- (8,0) -- (8,4) -- (0,4) --  cycle;
    \draw (2,2) circle (0.5cm);
    \draw (2,2) circle (1.25cm);
    \draw (6,2) circle (0.5cm);
    \draw (6,2) circle (1.25cm);
    \path node at (0.75,3.5) {\large $\Struct{U}$};
    \path node at (2,2) {\large $\Struct{M}$};
    \path node at (6,2) {\large $\Struct{N}$};
    \path node at (3.5,1) {\large $\Struct{M}^*$};
    \path node at (7.5,1) {\large $\Struct{N}^*$};
    \node (Idom) at (2.8,2) {};
    \node (Irng) at (5.2,2) {};
    \draw[<->, bend left] (Idom) to node {\large $\OLId{latin-looped}{I}$} (Irng) ;
  \end{tikzpicture}
  \caption{الـ!!{structure}~$\Struct{U}$ والتشاكل الجزئي الداخلي.}
\end{figure}

وموضع الحجة أن !!{language} الـ!!{structure}~$\Struct{U}$
تجيز !!{sentence}~$!D_\OLMathNumeral{1}$ من المنطق المجرد~$\OLId{latin-looped}{L}$، صادقة في $\Struct{U}$، تعبّر عن أن
$\Struct{M} \models_\OLId{latin-looped}{L} !E$ و$\Struct{N} \not\models_\OLId{latin-looped}{L} !E$ (وهنا نستعمل
خاصية النسبنة). وتجيز كذلك !!{sentence}~$!D_\OLMathNumeral{2}$ \emph{من الرتبة الأولى}
تصدق في $\Struct{U}$ وتقول إن $\Struct{M} \simeq_\OLId{latin-ordinary}{p} \Struct{N}$ بواسطة
التشاكل الجزئي~$\OLId{latin-looped}{I}$. فتقتضي خاصية لوفنهايم--سكولم صدق $!D_\OLMathNumeral{1}$ و$!D_\OLMathNumeral{2}$ معًا
في نموذج !!a{enumerable}~$\Struct{U_0}$ يحتوي بنيتين جزئيتين متشاكلتين
جزئيًا $\Struct{M_0}$ و$\Struct{N_0}$ بحيث
$\Struct{M_0} \models_\OLId{latin-looped}{L} !E$ و$\Struct{N_0} \not\models_\OLId{latin-looped}{L} !E$. وأما
الـ!!{structure}s الجزئية !!{enumerable} فمتى تشاكلت جزئيًا لزم تشاكلها
بحسب \olref[bas][pis]{thm:p-isom1}. وهذا يناقض خاصية التشاكل في المنطقات
المجردة السوية، فبطل الفرض.
\end{proof}

\end{document}
